% page 319 of the facsimile (image p-318 of the scan)
% block 165.1 x 267.9 mm ; left margin 36.9 mm
\pgc{15.240mm}{0.000mm}{
\l{}
\l{}
\l{}
\l{}
\l{\cel{52}311}
\l{}
\l{\sou{kovrar}.\cel{2}(\sou{trans., per,}\cel{1}ye). I. (1) Garnisar ye ulo quan onu}
\l{\cel{3}pozas sur...; (2) Garnisar per esar pozita sur...- II.}
\l{\cel{3}(1) Dissemar ulo sur ulo, disjete ; (2) Dissemesir disjete}
\l{\cel{3}sur ulo. - III. (1) Protektar per ulo quan onu pozas sur...;}
\l{\cel{3}(2) Protektar per esar sur, avan. - IV. (1) Celar per pozar}
\l{\cel{3}ulo sur, avan...; (2) Celar per esar sur, avan...- EFIS.}
\l{}
\l{\sou{koxalgio.}\cel{1}(\sou{patol.)}\cel{3}Artrito o tumoro blanka ye la hancho.- DEF.}
\l{}
\l{\sou{kozo}. Omna realajo, konkreta od abstraktita, quan onu indikas}
\l{\cel{3}ye maniero nedeterminita, nepreciza. -\cel{2}FIS.}
\l{}
\l{kozako.Slavo, kelke mestica, qua habitas regioni di la Rusia}
\l{\cel{3}di Europa e di Azia ; kom kavalriano ecelanta, lu posedis}
\l{\cel{3}privileji administreriala e fiskala, kondicionata per lua}
\l{\cel{3}devo esar armeano dum duadek yari, e, pose, en la milico}
\l{\cel{3}til la fino di sua valideso.- DEFIRS.}
\l{}
\l{\sou{krabo}. (\sou{zool.}) Nomo vulgara di krustacei diversa, dekpeda,}
\l{\cel{3}kurta-kauda, di qua la tipo esas la krabo ordinara (L.\cel{1}\sou{cancer,}}
\l{\cel{3}\sou{carcinus}), speco manjebla qua habitas en la sablo di maro.}
\l{\cel{3}- DEFR.}
\l{}
\l{\sou{krago.}\cel{2}Peco de telo, de dentelo, qua, pos cirkumirir la kolo,}
\l{\cel{3}retrovenas ad la supra parto di la pektoro. - D.}
\l{}
\l{\sou{krak}\cel{1}!\cel{2}Interjeciono qua expresas evento subita, o la bruiso de}
\l{\cel{3}krako.}
\l{}
\l{krakar. (netrans.) Produktar bruiso sika, per la krumplo di la}
\l{\cel{3}parti. - DEF.}
\l{}
\l{\sou{kraknelo.}\cel{2}Kukajo harda, qua krakas sub la denti qui mastikas}
\l{\cel{3}lu. - EF.}
\l{}
\l{\sou{krampo}. (\sou{patol.})\cel{2}Kontrakto dolorigiva qua produktas rigideso}
\l{\cel{3}kurta-dura di la muskuli brakiala, gambala, (\sou{analogie}\cel{1}:}
\l{\cel{3}stomakala). Kontrakto spasmala. - DEF.}
\l{}
\l{\sou{krampono}. La signo tipografiala :}
\l{}
\l{\sou{krano.}\cel{2}(\sou{tekn.})\cel{2}Aparato por sublevar la kargaji e transportar}
\l{\cel{3}li kurta-diste. Lu konsistas ek framo ek fero od ek ligno,}
\l{\cel{3}muntita sur pivoto cirkum qua lu povas rotacar, ed ek du}
\l{\cel{3}mekanismi : un mekanismo di sublevo (puliaro o vindilo) ed}
\l{\cel{3}un mekanismo di rotacigo. - DER.}
\l{}
\l{kranio. (\sou{anat.)}\cel{2}Buxo osta qua kontenas e protektas la cerebro.-}
\l{\cel{3}- EFIS.}
\l{}
\l{krapular. (netrans.) Debochar grosiere ; karakterizivo :}
\l{\cel{3}eceso di vino-drinko. - DEFIS.}
\l{}
\l{kraso. Strato de materio sordida qua formacesas ed akumulesas}
\l{\cel{4}sur la pelo, sur linjo, sur la vesti quin onu +weras, sur la}
\l{\cel{3}objekti. - FS.}
}
